% Part: sets-functions-relations
% Chapter: relations
% Section: relations-as-sets

\documentclass[../../../include/open-logic-section]{subfiles}

\begin{document}

\olfileid{sfr}{rel}{set}
\olsection{Relasi minangka Himpunan}

\begin{explain}
Ing \olref[sfr][set][imp]{sec}, kita nyebut sawetara himpunan
wigati: $\Nat$, $\Int$, $\Rat$, $\Real$. Mesthine kita eling
sawetara relasi kang narik kawigaten ing antarane !!{element}s
saka sawetara himpunan iki. Tuladhane, saben himpunan kasebut
duwe \emph{relasi tatanan} kang wis lumrah banget.
Uga ana relasi \emph{padha karo}, kang mung diduweni saben objek
marang awake dhewe lan ora marang liyane.
Ana luwih akeh relasi kang narik kawigaten kang bakal kita temoni,
lan luwih akeh maneh relasi kang bisa ana. Nanging sadurunge
ngrembug iku, kita bakal miwiti kanthi nuduhake yen relasi bisa
dianggep minangka jinis himpunan khusus.

Kanggo iki, elinga rong bab saka \olref[sfr][set][pai]{sec}.
Kapisan, elinga gagasan \emph{pasangan urut}: kanthi ana $a$ lan $b$,
kita bisa mbentuk~$\tuple{a, b}$. Ing kene urutan anggota
\emph{pancen} wigati. Mula yen $a \neq b$,
banjur $\tuple{a, b} \neq \tuple{b, a}$.
(Bandhingna karo pasangan tanpa urutan, yaiku himpunan kang
anggotane $2$, kanthi $\{a, b\}=\{b, a\}$.)
Kapindho, elinga gagasan \emph{produk Kartesius}:
yen $A$ lan $B$ iku himpunan, kita bisa mbentuk~$A \times B$,
himpunan kabeh pasangan $\tuple{x, y}$ kanthi $x \in A$
lan $y \in B$. Mligine, $A^{2}= A \times A$ iku
himpunan kabeh pasangan urut saka~$A$.

Saiki kita ngrembug relasi tartamtu ing sawijining himpunan:
relasi-$<$ ing himpunan~$\Nat$ saka wilangan natural.
Gatekna himpunan kabeh pasangan wilangan $\tuple{n, m}$
kanthi $n<m$, yaiku,
\[
  R=\Setabs{\tuple{n, m}}{n, m \in \Nat \text{ lan } n<m}.
\]
Ana gayutan kang raket antarane $n$ kurang saka $m$ lan
pasangan $\tuple{n, m}$ dadi anggota $R$, yaiku:
\[
      n<m\text{ yen lan mung yen }\tuple{n, m} \in R.
\]
Tanpa kelangan informasi, kita bisa nganggep himpunan $R$
\emph{minangka} relasi-$<$ ing $\Nat$.

Kanthi cara kang padha, kita bisa mbentuk subhimpunan saka
$\Nat^{2}$ kanggo saben relasi antarane wilangan.
Kosok baline, kanggo saben himpunan pasangan wilangan
$S \subseteq \Nat^{2}$, ana relasi wilangan kang cocog:
yaiku relasi kang diduweni $n$ marang $m$ yen lan mung yen
$\tuple{n, m} \in S$. Iki mbenerake definisi ing ngisor iki:
\end{explain}

\begin{defn}[Relasi biner]
\emph{Relasi biner} ing himpunan $A$ yaiku subhimpunan
saka~$A^{2}$. Yen $R\subseteq A^{2}$ iku relasi biner
ing~$A$ lan $x, y \in A$, kadhangkala kita nulis
$Rxy$ (utawa $xRy$) kanggo $\tuple{x, y} \in R$.
\end{defn}

\begin{ex}
  \ollabel{relations}
Himpunan $\Nat^{2}$ saka pasangan wilangan natural bisa
ditata ing matriks rong dimensi kaya mangkene:
\[
  \begin{array}{ccccc}
  \mathbf{\tuple{ 0,0 }} & \tuple{ 0,1 } &
    \tuple{ 0,2 } & \tuple{ 0,3 } & \ldots\\
  \tuple{ 1,0 } & \mathbf{\tuple{ 1,1 }} &
    \tuple{ 1,2 } & \tuple{ 1,3 } & \ldots\\
  \tuple{ 2,0 } & \tuple{ 2,1 } &
    \mathbf{\tuple{ 2,2 }} & \tuple{ 2,3 } & \ldots\\
  \tuple{ 3,0 } & \tuple{ 3,1 } & \tuple{ 3,2 } &
    \mathbf{\tuple{ 3,3 }} & \ldots\\
  \vdots & \vdots & \vdots & \vdots & \mathbf{\ddots}
  \end{array}
\]
Diagonal ing kene dicithak kandel, amarga subhimpunan saka
$\Nat^2$ kang ngemot pasangan ing diagonal, yaiku,
\[
  \{\tuple{0,0 }, \tuple{ 1,1 }, \tuple{ 2,2 }, \dots\},
  \]
iku \emph{relasi identitas ing}~$\Nat$.
(Amarga relasi identitas kerep digunakake, kita netepake
$\Id{A}=\Setabs{\tuple{ x,x }}{x \in A}$ kanggo saben himpunan $A$.)
Subhimpunan kabeh pasangan kang ana ing sadhuwure diagonal, yaiku,
\[
  L = \{\tuple{ 0,1 },\tuple{ 0,2 },\ldots,\tuple{ 1,2 },
  \tuple{ 1,3 }, \dots, \tuple{ 2,3 }, \tuple{ 2,4 },\ldots\},
\]
iku relasi \emph{kurang saka}, yaiku $Lnm$ yen lan mung yen $n<m$.
Subhimpunan pasangan kang ana ing sangisore diagonal, yaiku,
\[
  G=\{ \tuple{ 1,0 },\tuple{ 2,0 },\tuple{
    2,1 }, \tuple{ 3,0 },\tuple{ 3,1 },\tuple{ 3,2 }, \dots\},
\]
iku relasi \emph{luwih saka}, yaiku $Gnm$ yen lan mung yen $n>m$.
Gabungan $L$ karo $I$, kang bisa kita jenengi $K=L\cup I$,
iku relasi \emph{kurang saka utawa padha karo}:
$Knm$ yen lan mung yen $n \le m$.
Semono uga, $H=G \cup I$ iku \emph{relasi luwih saka
utawa padha karo}. Relasi $L$, $G$, $K$, lan $H$ iki
jinis relasi khusus kang diarani \emph{tatanan}.
$L$ lan $G$ duwe sipat yen ora ana wilangan kang nduweni
relasi $L$ utawa $G$ marang awake dhewe
(yaiku, kanggo kabeh $n$, ora $Lnn$ lan ora $Gnn$).
Relasi kang duwe sipat iki diarani \emph{irrefleksif};
yen relasi kasebut uga tatanan, diarani \emph{tatanan ketat}.
\end{ex}

\begin{explain}
Sanajan tatanan lan identitas iku relasi kang wigati lan lumrah,
kudu ditegesake yen miturut definisi kita, \emph{saben}
subhimpunan saka $A^{2}$ iku relasi ing~$A$,
ora preduli sepira ora lumrah utawa katon digawe-gawe.
Mligine, $\emptyset$ iku relasi ing saben himpunan
(\emph{relasi kosong}, kang ora diduweni pasangan anggota apa wae),
lan $A^{2}$~dhewe uga relasi ing~$A$
(kang diduweni saben pasangan), diarani \emph{relasi universal}.
Nanging bab kaya
$E=\Setabs{\tuple{n, m}}{n>5 \text{ utawa } m \times n \ge 34}$
uga dianggep relasi.
\end{explain}

\begin{prob}
  Dhaptarna !!{element}s relasi $\subseteq$ ing himpunan
  $\Pow{\{a, b, c\}}$.
\end{prob}

\end{document}
